% 克莱因瓶示意 (Klein Bottle)
% 描述: 展示克莱因瓶的构造，强调其无边界且不可定向的性质
% 数学概念: 克莱因瓶、闭曲面、不可定向、商空间、浸入
\documentclass[border=10pt]{standalone}
\usepackage{tikz}
\usepackage{amsmath,amssymb}
\usetikzlibrary{arrows.meta,positioning,calc,patterns,decorations.pathmorphing}

\begin{document}
\begin{tikzpicture}[
    >=Stealth,
    scale=1.0
]

% ============= 左侧: 二维商空间表示 =============
\begin{scope}[xshift=-5cm, yshift=0.5cm]
    % 标题
    \node[font=\bfseries] at (0, 3) {商空间构造};
    
    % 矩形
    \fill[cyan!15] (-2, -1.5) rectangle (2, 1.5);
    \draw[line width=1.5pt, color=blue!60!black] (-2, -1.5) rectangle (2, 1.5);
    
    % 水平边 (反向)
    \draw[color=green!60!black, line width=2pt, ->] (-1.8, 1.5) -- (-0.2, 1.5);
    \draw[color=green!60!black, line width=2pt, ->] (0.2, 1.5) -- (1.8, 1.5);
    \draw[color=green!60!black, line width=2pt, ->] (-1.8, -1.5) -- (-0.2, -1.5);
    \draw[color=green!60!black, line width=2pt, ->] (0.2, -1.5) -- (1.8, -1.5);
    \node[color=green!60!black, above] at (0, 1.5) {$a$};
    \node[color=green!60!black, below] at (0, -1.5) {$a^{-1}$};
    
    % 垂直边 (反向)
    \draw[color=red!70!black, line width=2pt, ->] (-2, -1.3) -- (-2, 1.3);
    \draw[color=red!70!black, line width=2pt, ->] (2, 1.3) -- (2, -1.3);
    \node[color=red!70!black, left] at (-2, 0) {$b$};
    \node[color=red!70!black, right] at (2, 0) {$b^{-1}$};
    
    % 扭曲标记
    \node[font=\scriptsize, color=orange!80!black] at (0, 0) {\textbf{两边都反向粘合}};
    
    % 公式
    \node[anchor=north, font=\small] at (0, -2.2) {$K = I \times I / \sim$};
    \node[anchor=north, font=\scriptsize, text width=5cm, align=center] 
        at (0, -2.8) {$(0, t) \sim (1, t)$, $(s, 0) \sim (1-s, 1)$};
\end{scope}

% ============= 中间: "8字形"浸入表示 =============
\begin{scope}[xshift=0.5cm, yshift=0.5cm]
    % 标题
    \node[font=\bfseries] at (0, 3) {浸入 $\mathbb{R}^3$};
    
    % 克莱因瓶的简化表示 (类似两个莫比乌斯带粘合)
    % 底部大圆
    \draw[color=blue!60!black, line width=1.5pt] 
        (0, -1.5) ellipse (1.5 and 0.4);
    
    % 瓶身扭曲上升
    \draw[color=blue!60!black, line width=1pt, fill=cyan!15, opacity=0.7]
        (-1.5, -1.5) .. controls (-1.8, -0.5) and (-1.5, 0.5) .. (-0.8, 0.8)
                     .. controls (0, 1.1) and (0.8, 0.8) .. (1.5, 0.5)
                     .. controls (1.8, 0.2) and (1.5, -0.5) .. (1.5, -1.5);
    
    % 瓶颈穿过自身
    \draw[color=blue!60!black, line width=1pt, dashed, opacity=0.5]
        (-0.3, 0.3) .. controls (0, 0.5) .. (0.3, 0.3);
    
    % 顶部小圆 (瓶口)
    \draw[color=blue!60!black, line width=1.5pt] 
        (0, 0.8) ellipse (0.8 and 0.25);
    
    % 指示箭头
    \draw[->, color=orange!80!black, line width=1.2pt] 
        (1.8, -0.5) arc (-30:30:1 and 0.8);
    \node[color=orange!80!black, font=\scriptsize, right] at (2.2, -0.2) {自交};
    
    % 说明
    \node[anchor=north, font=\scriptsize, text width=5cm, align=center] 
        at (0, -2.5) {无法无自交地嵌入$\mathbb{R}^3$，\\
                       但可以嵌入$\mathbb{R}^4$};
\end{scope}

% ============= 右侧: 两个莫比乌斯带 =============
\begin{scope}[xshift=5.5cm, yshift=0.5cm]
    % 标题
    \node[font=\bfseries] at (0, 3) {分解结构};
    
    % 两个莫比乌斯带示意图
    % 第一个莫比乌斯带
    \fill[purple!20, opacity=0.6] (-1.5, 0.5) ellipse (1.2 and 0.6);
    \draw[color=purple!70!black, line width=1pt] (-1.5, 0.5) ellipse (1.2 and 0.6);
    \node[font=\scriptsize, color=purple!70!black] at (-1.5, 0.5) {$M_1$};
    
    % 第二个莫比乌斯带
    \fill[orange!20, opacity=0.6] (0.8, 0.5) ellipse (1.2 and 0.6);
    \draw[color=orange!70!black, line width=1pt] (0.8, 0.5) ellipse (1.2 and 0.6);
    \node[font=\scriptsize, color=orange!70!black] at (0.8, 0.5) {$M_2$};
    
    % 粘合指示
    \draw[<->, line width=1.5pt, color=red!70!black] 
        (-0.3, 0.5) -- (0.2, 0.5);
    \node[color=red!70!black, font=\scriptsize, above] at (-0.05, 0.6) {沿边界粘合};
    
    % 等式
    \node[anchor=north, font=\small] at (0, -0.8) {$K = M_1 \cup_{\partial} M_2$};
    
    % 说明
    \node[anchor=north, font=\scriptsize, text width=5cm, align=center] 
        at (0, -1.5) {克莱因瓶是两个莫比乌斯带\\沿边界粘合而成};
\end{scope}

% ============= 底部: 数学性质 =============
\node[anchor=north, font=\small, draw=cyan!30, fill=cyan!5, rounded corners, inner sep=10pt] 
    at (4, -4.5) {
    \begin{tabular}{ll}
        \textbf{拓扑性质:} & \\
        边界: & $\partial K = \emptyset$ (闭曲面) \\
        可定向性: & 不可定向 \\
        亏格: & 不可定向亏格 2 \\
        Euler示性数: & $\chi(K) = 0$ \\
        基本群: & $\pi_1(K) = \langle a, b \mid aba^{-1}b = 1 \rangle$ \\
        分类: & 不可定向闭曲面
    \end{tabular}
};

% ============= 基本群图示 =============
\begin{scope}[xshift=-5cm, yshift=-4.5cm]
    \node[anchor=north west, font=\scriptsize, text width=5cm, align=left] 
        at (0, 0) {
        \textbf{基本群:} \\
        $\pi_1(K) = \langle a, b \mid abab^{-1} = 1 \rangle$\\[3pt]
        非阿贝尔群，\\
        包含挠元
    };
\end{scope}

\end{tikzpicture}
\end{document}
